\documentclass[runningheads]{llncs}

\usepackage{graphicx}
\usepackage[final,hidelinks]{hyperref}
\usepackage{listings}
\usepackage{xcolor}
%\usepackage{todonotes} %[disable]
\usepackage{enumitem}
\usepackage{lstcoq}
\lstset{language=coq}
\usepackage{proof-dashed}

\usepackage{todonotes}

\definecolor{dkviolet}{rgb}{.5,0,.5}
\definecolor{dkblue}{rgb}{0,0,.5}
\definecolor{dkgreen}{rgb}{0,.5,0}
\definecolor{dkred}{rgb}{.5,0,0}
\definecolor{ltblue}{rgb}{0,.4,.6}

\usepackage{comment}

% If you use the hyperref package, please uncomment the following line
% to display URLs in blue roman font according to Springer's eBook style:
% \renewcommand\UrlFont{\color{blue}\rmfamily}

\begin{document}
\title{Certifying Choreography Compilation\thanks{Work partially supported by Villum Fonden, grant no.\ 29518.}}

\author{Luís Cruz-Filipe\orcidID{0000-0002-7866-7484}
\and
Fabrizio Montesi\orcidID{0000-0003-4666-901X}
\and
Marco Peressotti\orcidID{0000-0002-0243-0480}
}

\authorrunning{L. Cruz-Filipe et al.}

\institute{Department of Mathematics and Computer Science, University of Southern Denmark\\
Campusvej 55, 5230 Odense M, Denmark
\email{\{lcfilipe,fmontesi,peressotti\}@imada.sdu.dk}}

\maketitle

\begin{abstract}
Choreographic programming is a paradigm for developing concurrent and distributed systems, where programs are \emph{choreographies} that define, from a global viewpoint, the computations and interactions that communicating processes should enact.
\emph{Choreography compilation} translates choreographies into the local definitions of process behaviours, given as terms in a process calculus.

Proving choreography compilation correct is challenging and error-prone, because it requires relating languages in different paradigms (global interactions vs local actions) and dealing with a combinatorial explosion of proof cases.
We present the first certified program for choreography compilation for a nontrivial choreographic language supporting recursion.
% In addition to the increase in confidence on the viability of the choreographic programming approach, our development, which we carry out using the Coq theorem prover, makes explicit several properties that are implicitly assumed in papers.
\keywords{Choreographic Programming \and Formalisation \and Compilation.}
\end{abstract}

\section{Introduction}\label{sec:intro}
Choreographic programming is an emerging programming paradigm, where the desired communication behaviour of a system of communicating processes can be defined from a global viewpoint in programs known as choreographies~\cite{M13p}.
Then, a provably-correct compiler can automatically generate executable code for each process, providing the guarantee that executing these processes together will implement the communications prescribed in the choreography~\cite{CHY12,CM13}. The theory of such compilers is typically called EndPoint Projection (EPP).

Choreographies are inspired by the ``Alice and Bob'' notation for security protocols~\cite{NS78}. The key idea is to have a linguistic primitive for a communication from a process (also called role, or participant) to another. For example, the statement \lstinline+Alice.e -> Bob.x+ reads ``\lstinline+Alice+ evaluates expression \lstinline+e+ and sends the result to \lstinline+Bob+, who stores it in variable \lstinline+x+''.
This syntax has two main advantages. The first is that it makes the desired communications syntactically manifest in a choreography, making choreographic programming suitable for making interaction protocols precise.
The second is that the syntax does not allow for writing mismatched send and receive actions by design, which makes code generated from a choreography enjoy progress (the system never gets stuck)~\cite{CM13}.

The potential of choreographic programming has motivated the study of choreographic languages and EPP definitions for different applications, including self-adaptive systems~\cite{DGGLM17}, information flow~\cite{LN15}, system integration~\cite{GLR18}, parallel algorithms~\cite{CM16}, cyber-physical systems~\cite{GMP20,LH17,LNN16}, and security protocols~\cite{GMP20}.

Proving a methodology for EPP correct involves three elements: the syntax and semantics of the source choreography language and of the target language, and the definition of the compiler. A single instruction at the choreographic level might be implemented by multiple instructions in the target language. All this makes formulating a theory of choreographic programming error-prone: for even simpler approaches, like abstract choreographies without computation, it has been recently discovered that a few key results published in peer-reviewed articles do not hold and their theories required adjustments~\cite{SY19}, raising concerns about the soundness of these methods.

In this article, we present a certified program for EPP, which translates terms of a Turing complete choreographic language into terms of a distributed process calculus. Our main result is the formalisation of the hallmark result of choreographic programming, the ``EPP Theorem'': an operational correspondence between choreographies and their endpoint projections.
This is the first time that such a result has been formalised in a theorem prover, increasing our confidence in the methodology of choreographies.

\paragraph{Structure and Related Work.}
We use Coq for our development, assuming some familiarity with it. We start from a previous formalisation~\cite{CMP21} of a choreographic language (Core Choreographies~\cite{CM20}), which we recap in Sect.~\ref{sec:background}.

In Sect.~\ref{sec:sp}, we formalise our target language: a distributed process calculus inspired by the informal presentation in~\cite{M21itc}. The calculus has communication primitives that recall those commonly used for implementing choreography languages, e.g., as found in (multiparty) session types~\cite{HYC16}.

In Sect.~\ref{sec:merging}, we define \emph{merging}~\cite{CHY12}: a partial operator that addresses the standard problem (for choreographies) of checking that each process implementation eventually agrees on the choice between alternative behaviours in protocols~\cite{BBO12,CHY12}.

Building on merging, in Sect.~\ref{sec:epp}, we define EPP. Then, in Sect.~\ref{sec:pruning}, we explore \emph{pruning}~\cite{CHY12}: a preorder induced by merging that plays a key role in the EPP Theorem, which is proved in Sect.~\ref{sec:epp-theorem}.

Choreographies are used in industry for the specification and definition of web services and business processes~\cite{bpmn,wscdl}. These languages feature recursion or loops, which are not present in the only other formalisation work on choreographies that we are aware of~\cite{GA18}. We have validated the theory of EPP from~\cite{M21itc}, and made explicit several properties that are typically only implicitly assumed. Our results show that the ideas developed by researchers on choreographies, like merging, can be relied upon for languages of practical appeal.

\section{Background}
\label{sec:background}

We use the choreographic language of~\cite{M21itc}, which is inspired by Core Choreographies (CC)~\cite{CM20}. So far, this is the only Turing complete choreographic language that has been formalised~\cite{CMP21}.
%We use Core Choreographies (CC)~\cite{CM20} as choreographic language, in the variation presented in~\cite{M21itc}. As of now, this is the only available formalisation of a Turing complete choreographic language~\cite{CMP21}.
In this section, we recap this formalisation, which our work builds upon.
% , but we
% only include a minimum of Coq code.
We refer the reader to~\cite{CMP21} for a discussion of the design options both behind the
choreographic language and the formalisation.
Some of these are relevant for our development, and we explain them when needed.

The purpose of CC is to model communication protocols involving several participants, called processes.
Each process is equipped with memory cells, identified by variables.
Communications are of two kinds: value communications, where a process evaluates an expression that
may depend on the values stored in its memory and sends the result to a (distinct) process; and
label selections, where a process selects from different behaviours available at another process by means of an identifier (the label).
Label selections are typically used to communicate local choices made by a process to other processes.
In order to model recursive or infinite behaviour, CC supports the definition of procedures.
% , which may be called by their name.

\paragraph{Syntax.}
The formalisation of CC is parametric on the types of processes (\lstinline!Pid!), variables
(\lstinline!Var!), expressions (\lstinline!Expr!), values (\lstinline!Value!,
stored in memory), Boolean expressions (\lstinline!BExpr!)
% , for writing conditionals)
and procedure
names (\lstinline!RecVar!).
These are required to have decidable syntactic equality.
Labels are restricted to a two-element set (\lstinline!left! and
\lstinline!right!), a common choice in choreographies and session types~\cite{CP10,CLMSW16,CMS18,W14}.

The syntax of choreographies is defined by the following BNF grammar.\todo{we simplify the coq notation}
\begin{lstlisting}
Eta ::= p#e --> qdollarx | p --> q[l]
C ::= Eta;; C | If p ?? b Then Ct Else Ce | Call X | RT_Call X ps C | End
\end{lstlisting}
\lstinline!Eta! is the type of communication actions, where \lstinline+p#e --> qdollarx+ is a value communication and \lstinline+p --> q[l]+ is a label selection.
Choreography \lstinline!eta;; C!, can execute a communication \lstinline!eta! and continue as \lstinline!C!.
Conditionals \lstinline+If p ?? b Then Ct Else Ce+ evaluate of \lstinline+b+ and continue as either \lstinline+Ct+ or \lstinline+Ce+, according to whether \lstinline+b+ evaluates to \lstinline+true+ or \lstinline+false+.
\lstinline+Call X+ denotes a call to a procedure \lstinline+X+.
Since processes execute independently, the runtime term \lstinline+RT_Call X ps C+ is used to denote
the situation where execution of \lstinline+X+ has started and reduced to \lstinline+C+, but the
processes in \lstinline+ps+ still have not entered \lstinline+X+.
Term \lstinline+End+ is the terminated choreography.

This grammar is implemented as a Coq inductive type, e.g., \lstinline+eta;; C+ stands for \lstinline+Interaction eta C+, where \lstinline+Interaction : Eta -> Choreography -> Choreography+ is a constructor of type \lstinline+Choreography+.

%% The syntax of choreographies is defined in Coq by the following inductive types, which directly correspond to a BNF grammar.
%% \begin{lstlisting}
%% Inductive Eta : Type :=
%%  | Com : Pid -> Expr -> Pid -> Var -> Eta   | Sel : Pid -> Pid -> Label -> Eta.

%% Inductive Choreography : Type :=
%%  | Interaction : Eta -> Choreography -> Choreography
%%  | Cond : Pid -> BExpr -> Choreography -> Choreography -> Choreography
%%  | Call : RecVar -> Choreography
%%  | RT_Call : RecVar -> (list Pid) -> Choreography -> Choreography
%%  | End : Choreography.
%% \end{lstlisting}
%% \lstinline!Eta! is the type of communication actions, where \lstinline|Com| is a value communication and \lstinline|Sel| is a label selection. These come with the more suggestive notations \lstinline+p#e --> qdollarx+, for \lstinline+Com p e q x+, and \lstinline+p --> q[l]+, for \lstinline+Sel p q l+.

%% Choreography \lstinline!Interaction e C!, also written as \lstinline!e;; C!, can execute a
%% communication \lstinline!e! and continue as \lstinline!C!.
%% Conditionals are written \lstinline+Cond p b Ct Ce+, or \lstinline+If p ?? b Then Ct Else Ce+, where
%% after evaluation of \lstinline+b+ the protocol continues as either \lstinline+Ct+ or \lstinline+Ce+,
%% according to whether \lstinline+b+ evaluates to \lstinline+true+ or \lstinline+false+.
%% A call to a procedure \lstinline+X+ is written \lstinline+Call X+ (there is no continuation for procedure calls).
%% Since processes execute independently, the runtime term \lstinline+RT_Call X ps C+ is used to denote
%% the situation where execution of \lstinline+X+ has started and reduced to \lstinline+C+, but the
%% processes in \lstinline+ps+ still have not entered \lstinline+X+.
%% Term \lstinline+End+ is the terminated choreography.

Intuitively, \lstinline+Call X+ behaves as the definition of \lstinline+X+.
However, the semantics includes internal actions at each process \lstinline+p+, corresponding to
\lstinline+p+ calling the procedure.
Thus, when the first process \lstinline+p+ calls \lstinline+X+, choreography \lstinline+Call X+ transitions to \lstinline+RT_Call X ps C+, where \lstinline+C+ is the definition of \lstinline+X+ and
\lstinline+ps+ are the processes used in \lstinline+X+ (other than \lstinline+p+).
This term can then either execute by another process entering the procedure (and this process is
removed from \lstinline+ps+), or by executing some action in \lstinline+C+ that does not involve
processes in \lstinline+ps+ (and \lstinline+C+ is updated).
When the last process in \lstinline+ps+ enters the procedure call, the term transitions simply to
\lstinline+C+.

This informal explanation illustrates two important properties of CC.
First, the semantics is parametric on the set of procedure definitions, modelled as a function \lstinline+Defs : DefSet+, where \lstinline+DefSet = RecVar -> (list Pid)*Choreography+, where \lstinline+fst (Defs X)+ is the set of processes involved in \lstinline+X+, and \lstinline+snd (Defs X)+ is the choreography defining \lstinline+X+.
A \emph{program} is a pair consisting of a set of procedure definitions, \lstinline+Procedures P+, and a \lstinline+Choreography+, \lstinline+Main P+.
We write \lstinline+Vars P X+ for \lstinline+fst (Procedures P X)+ and \lstinline+Procs P X+
for \lstinline+snd (Procedures P X)+.

The second property is that CC supports out-of-order execution: interactions involving distinct processes can be executed in any order, reflecting concurrency.
% in the choreography can
% be executed, as long as the processes involved in it do not appear previously in the choreography.
For example, \lstinline+p#e --> qdollarx;; r --> s[left]+ can execute as a
communication from \lstinline+p+ to \lstinline+q+ followed by a label selection from \lstinline+r+
to \lstinline+s+, but also as the latter label selection followed by the former communication.
The semantics of CC is formalised in~\cite{CMP21}; this informal explanation suffices for the remainder of this article.

The runtime term \lstinline+RT_Call X ps C+ is not meant to be used when writing a choreography, and
in particular it should not occur in procedure definitions.
A choreography that does not include any such term is called \emph{initial}.

There are several restrictions when writing choreographies, of three kinds.
% , which fall into three
% categories.
\begin{enumerate}[nosep,noitemsep,label={(\roman{*})}]
\item Restrictions on the intended use of choreographies.
% related to the intended usage of choreographies.
  Interactions must have distinct processes (no self-communication), e.g., 
  \lstinline+p#e --> pdollarx+ is disallowed.
%   Both \lstinline+Main P+ and \lstinline+Procs P X+ should satisfy this condition.

\item Restrictions on the intended use of runtime terms.
  All choreographies \lstinline+Procs P X+ must be initial.
  As for \lstinline+Main P+, for any subterm \lstinline+RT_Call X ps C+,
  \lstinline+ps+ must be nonempty and include only process names that occur in
  \lstinline+Vars P X+.

\item Restrictions that arise from design choices in the formalisation.
  Informally, \lstinline+Vars X+ contains the processes that are used in \lstinline+Procs X+.
  This introduces constraints that are encapsulated in the definition of well-formed
  program.
\end{enumerate}
The constraints in the last category are particularly relevant for the proof of the EPP Theorem, so we discuss them in that context in Sect.~\ref{sec:epp-theorem}.
% It is easier to understand them in that context, so we do not discuss them further at this stage.

\paragraph{Semantics.}
The semantics of CC is defined as a labelled transition system using inductive types.
It uses a notion of \emph{state}, which is simply a function mapping process variables to their values: \lstinline+State:Pid -> Var -> Value+.

The semantics is structured in three layers.
The first layer specifies transitions with the following relation, parameterised on a
set of procedure definitions.
\begin{lstlisting}
CCC_To : DefSet -> Choreography -> State
				-> RichLabel -> Choreography -> State -> Prop
\end{lstlisting}
Rules in this set include that \lstinline+p#e --> qdollarx;; C+ in state \lstinline+s+
transitions to \lstinline+C+ in state \lstinline+s'+, where \lstinline+s'+ coincides with \lstinline+s+
except that \lstinline+s' q x+ now stores the value obtained by evaluating \lstinline+e+
at \lstinline+p+ in state \lstinline+s+.\footnote{The semantics of CC only requires extensional equality of states, which is why \lstinline+s'+ is quantified over rather than directly defined from \lstinline+s+.}
The label in the transition includes information about the executed term -- the above
communication is labelled \lstinline+R_Com p v q x+.

% Out-of-order execution is modelled by rules that allow prefixing choreographies in reductions.
% For example, if \lstinline+CCC_To Defs C s t C' s'+ and \lstinline+eta+ is a communication action
% not involving any processes in \lstinline+t+, then also
% \lstinline+CCC_To Defs (eta;; C) s t (eta;; C') s'+.

The second layer raises transitions to the level of programs, abstracting from 
unobservable details.
It is defined by a single rule: if \lstinline+CCC_To Defs C s t C' s'+, then
\lstinline+({|Defs; C|},s) --[forget t]--> ({|Defs; C'|},s')+.
Here, \lstinline+{|Defs; C|}+ denotes\footnote{This is a slightly simplified version of the actual Coq notation.} the program built from \lstinline+Defs+ and \lstinline+C+, and
\lstinline+forget+ removes unobservable details from transition labels (e.g.,
\lstinline+forget (R_Com p v q x)+ is \lstinline+L_Com p v q+).
Labels corresponding to conditionals and procedure calls all simplify to
\lstinline+L_Tau p+, denoting an internal action at \lstinline+p+.
% 
The third layer defines the reflexive and transitive closure of program transitions.

Important properties of CC include deadlock-freedom by design (any program \lstinline+P+ such that
\lstinline+Main P <> End+ reduces), confluence (if \lstinline+P+ can execute two different
sequences actions due to out-of-order execution, then the resulting programs can always reduce to a
common program and state), and Turing completeness.
% In particular, this implies a form of determinism: if executing a program terminates, then it does
% so in a uniquely defined state.
All these results are formalised in~\cite{CMP21}.

% The formalisation also includes the result that CC is Turing-complete, proven by a straightforward
% adaptation of the proof in~\cite{CM20}.
% In (the remainder of) this article, we define a process calculus in which CC can be implemented,
% identify conditions for a choreography to be implementable, define a mechanical translation from such
% choreographies to implementations, and prove the translation correct.
% that choreographies are operationally correspondent to their generated implementations.

\todo{example of transitions of the previous example}

\section{Stateful Processes}
\label{sec:sp}

We now move to the contributions of this article. In this section, we formalise SP (Stateful Processes), the process calculus for implementing choreographies in~\cite{M21itc}, which is inspired by the process calculi in~\cite{CM17,CM20}.

\subsection{Syntax}
SP is defined as a Coq functor with the same parameters as CC.
Its syntax contains three main ingredients: behaviours, which define the actions performed by the individual processes; networks, which consist of several processes running in parallel; and
programs, which (as in CC) make available a common set of procedure definitions that each process
can use.
%\footnote{For programs obtained from choreographies, we later show that the sets of
% procedures that each process needs to access are disjoint, so this set can be distributed.}

\paragraph{Behaviours.}
Behaviours are sequences of local actions, corresponding to the possible terms that can be written in CC,
% built from the local actions corresponding to the possible interactions that can be
% written in CC.
% Value communications and local selections give two local actions each, while conditionals and
% procedure calls give one local action.
defined by the following BNF grammar.
\begin{lstlisting}
B ::= End | p!e;B | p?x;B | p(+)l;B | p&mBl//mBr | If b Then Bt Else Be | Call X
\end{lstlisting}
This grammar is again formalised as a Coq inductive type \lstinline+Behaviour+.
\lstinline+End+ denotes a terminated behaviour.
Value communications are implemented by \lstinline+p!e; B+, which evaluates \lstinline+e+, sends the result to \lstinline+p+, and continues as \lstinline+B+, and by \lstinline+p?x; B+, which receives a value from \lstinline+p+, storing it at \lstinline+x+, and continues as \lstinline+B+.\footnote{Processes
communicate by name. In practice, these can be either process identifiers (cf.\ actors), network addresses, or session correlation data.}

Label selections are divided into \lstinline!p(+)l; B! -- sending the label \lstinline+l+ to \lstinline+p+ and continuing as \lstinline+B+ -- and \lstinline+p & mBl // mBr+, where one of \lstinline+mBl+ or
\lstinline+mBr+ is chosen according to the label selected by \lstinline+p+.
Both \lstinline+mBl+ and \lstinline+mBr+ have type \lstinline+option Behaviour+: a process does not need to offer behaviours corresponding to all possible labels.
Informally, branching terms are partial functions from labels to behaviours.
Our definition capitalises on the fact that there are only two labels to simplify the formalisation.

A conditional \lstinline+If b Then Bt Else Be+ evaluates the Boolean expression \lstinline+b+ to decide whether to continue as \lstinline+Bt+ or \lstinline+Be+.
Finally, \lstinline+Call X+ simply continues as the behaviour defining procedure \lstinline+X+.

%% The inductive definition of this type is written as follows in Coq.
%% \begin{lstlisting}
%% Inductive Behaviour : Type :=
%% | End : Behaviour
%% | Send : Pid -> Expr -> Behaviour -> Behaviour
%% | Recv : Pid -> Var -> Behaviour -> Behaviour
%% | Sel : Pid -> Label -> Behaviour -> Behaviour
%% | Branching : Pid -> option Behaviour -> option Behaviour -> Behaviour
%% | Cond : BExpr -> Behaviour -> Behaviour -> Behaviour
%% | Call : RecVar -> Behaviour.
%% \end{lstlisting}
%% % As in CC, 
%% \lstinline+End+ denotes a terminated behaviour.
%% Value communications are implemented by \lstinline+Send+ and \lstinline+Recv+.
%% Behaviour \lstinline+Send p e B+, also written \lstinline+p!e; B+, evaluates
%% \lstinline+e+, sends the result to \lstinline+p+, and continues as \lstinline+B+.
%% The dual behaviour \lstinline+Recv p x B+, or \lstinline+p?x; B+, corresponds to receiving a value
%% from \lstinline+p+, storing it at \lstinline+x+, and continuing as \lstinline+B+.\footnote{Processes
%% communicate by name. In practice, these can be either process identifiers (cf.\ actors), network addresses, or session correlation data.}

%% Label selections are divided into \lstinline+Sel p l B+, or \lstinline!p(+)l; B! -- sending the
%% label \lstinline+l+ to \lstinline+p+ and continuing as \lstinline+B+ -- and
%% \lstinline+Branching p Bl Br+, or \lstinline+p & Bl // Br+, where one of \lstinline+Bl+ or
%% \lstinline+Br+ is chosen according to the label selected by \lstinline+p+.
%% Both \lstinline+Bl+ and \lstinline+Br+ have type \lstinline+option Behaviour+, which is the Coq
%% version of the \verb!Maybe! monad: a process does not need to offer behaviours corresponding to all
%% possible labels.

%% A conditional \lstinline+Cond b Bt Be+, suggestively written \lstinline+If b Then Bt Else Be+,
%% evaluates the Boolean expression \lstinline+b+ to decide whether to continue as \lstinline+Bt+ or
%% \lstinline+Be+.
%% Finally, the term \lstinline+Call X+ simply continues as the behaviour defining procedure
%% \lstinline+X+.

% \paragraph{Design choices.}
%\todo{hard cut here}
\begin{comment}
\begin{remark}
Informally, branching terms are described as partial functions from labels to behaviours.
Since Coq only allows for the definition of total functions, there is some design freedom in
implementing these terms.
Our first attempt was to define
\lstinline+Branching : Pid -> (Label -> option Behaviour) -> Behaviour+, which more precisely
captures the informal intuition.
However, this turned out to be extremely problematic: since equality of functions is not
extensional, equality of behaviours becomes undecidable, and the automatically generated induction
principles for \lstinline+Behaviour+ were too weak.
Furthermore, many results about behaviours only held up to functional extensionality, which required
equipping \lstinline+Behaviour+ with a setoid structure -- significantly complicating the
formalisation.
Since the set of labels is finite -- an essential assumption in many results, which again was
not trivial to combine with this definition of behaviour -- we took the option of requiring it to
consist of exactly two elements, as is common practice in many other works on choreographies and session types~\cite{CP10,CLMSW16,CMS18,W14}.
This choice also makes our notation closer to what is commonly used in informal works.
\end{remark}
\end{comment}

The automatically generated induction principles for \lstinline+Behaviour+ do not allow us to inspect the subterms of branching terms. Therefore, we prove a stronger result that allows us to apply the induction hypothesis to these terms.

%% Since the constructor for branching terms uses functions on \lstinline+Behaviour+, the automatically
%% generated induction principles for \lstinline+Behaviour+ are not as strong as needed.
%% We follow the standard approach of defining the depth of a \lstinline+Behaviour+ as the depth of its
%% AST, and proving the following more robust induction principle by induction on this
%% depth.\footnote{Coq's type inference mechanism allows us to omit types of most universally
%%   quantified variables. We abuse this possibility to lighten the presentation, and use variable
%%   names that are suggestive of their type.}
%% \begin{lstlisting}
%% Theorem Behaviour_ind' : forall P:Behaviour -> Prop, P End ->
%%     (forall p e B, P B -> P (p!e; B)) ->  (forall p v B, P B -> P (p?v; B)) ->
%%     (forall p l B, P B -> P (p(+)l; B)) ->
%%     (forall p mB mB', (forall B, mB = Some B -> P B) ->
%%                       (forall B, mB' = Some B -> P B) -> P (p & mB // mB')) ->
%%     (forall b B1 B2, P B1 -> P B2 -> P (If b Then B1 Else B2)) -> (forall X, P (Call X)) ->
%%     forall B, P B.
%% \end{lstlisting}
%% The interesting clause is the one for branching terms, which could not be derived automatically by
%% Coq.

\paragraph{Networks.}
A network is a function with type \lstinline+Pid -> Behaviour+.
We define extensional equality for networks, written \lstinline+N == N'+, and prove that it is an
equivalence relation.
In practice, networks only use a finite number of processes, and are written as (finite) parallel
compositions of behaviours.
We introduce these particular constructions: \lstinline+N|N'+ is the parallel composition of \lstinline+N+ and \lstinline+N'+; \lstinline+p[B]+ is the network where \lstinline+p+ is mapped to \lstinline+B+ and all other processes to \lstinline+End+; and \lstinline+N ~~ p+ denotes network \lstinline+N+ where \lstinline+p+ is now mapped to \lstinline+End+.
%We prove that these networks behave as expected.
%, in particular that \lstinline+|+ is commutative when its arguments have disjoint support.
%% \begin{lstlisting}
%% Definition EmptyNet : Network := fun _ => End.
%% Definition Process (p:Pid) (B:Behaviour) : Network :=
%%   fun p' => if (Pid_dec p' p) then B else End.
%% Definition Par (N N':Network) :=
%%   fun p => if (Behaviour_eq_End_dec (N p)) then N' p else N p.
%% \end{lstlisting}
%% Here, \lstinline+Pid_dec+ is the term stating that equality of process names is decidable, and
%% \lstinline+Behaviour_eq_End_dec+ is a specialisation of decidability of behaviour equality
%% stating that we can decide whether a behaviour is equal to \lstinline+End+.
%% Parallel composition of networks is written \lstinline+N1 | N2+, while \lstinline+Network_rm N p+
%% is written \lstinline+N ~~ p+.
The following are examples of properties useful for proving results such as confluence.\footnote{Coq's type inference mechanism allows us to omit types of most universally quantified variables. We abuse this possibility to lighten the presentation, and use variable names that are suggestive of their type.}
\begin{lstlisting}
Lemma Network_rm_out : forall N p p', p <> p' -> (N ~~ p) p' = N p'.
Lemma Network_eq_cross : forall N N1 N2 p q r s Bp Bq Br Bs,
  p <> q -> p <> r -> p <> s -> q <> r -> q <> s -> r <> s ->
  N1 == N ~~ p ~~ q | p [Bp] | q [Bq] -> N2 == N ~~ r ~~ s | r [Br] | s [Bs] ->
  N2 ~~ p ~~ q | p [Bp] | q [Bq] == N1 ~~ r ~~ s | r [Br] | s [Bs].
\end{lstlisting}
Most of these results do not hold for (intensional) equality of \lstinline+Network+s.

\paragraph{Programs.}
Finally, a \lstinline+Program+ is a pair consisting of a set of procedure definitions (of type \lstinline+DefSetB = RecVar -> Behaviour+) and a \lstinline+Network+.

Programs, networks and behaviours should satisfy well-formedness properties similar to those of CC.
However, these properties are automatically ensured when networks are automatically generated from
choreographies, so we do not discuss them here.
They are included in the formalisation.

\todo{example network that implements the choreography example}

\subsection{Semantics}

The semantics of SP is defined by a labelled transition system.
Transitions for communications match dual actions in two processes, while conditionals and procedure
calls simply run locally.
We report the representative cases -- the complete definition can be found in the source formalisation~\cite{CMP21-source}.

\todo{introduce an == operator as syntactic sugar for eq\_state\_ext}

\todo{these rules in inference rule notation first, then show only S\_Com in the Coq code}

\[
\infer[\mbox{\lstinline+S_Com+}]{
	\mbox{\lstinline+(|Defs; N|), s --[L_Com p v q]--> (|Defs; N'|), s'+}
}{
	\begin{array}{c}
	\mbox{\lstinline+v := (eval_on_state e s p)+}
	\quad
    \mbox{\lstinline+N p = q!e; B+}
    \quad
    \mbox{\lstinline+N q = p?x; B'+}
    \\
    \mbox{\lstinline+N' == N ~~ p ~~ q | p[B] | q[B']+}
    \quad
    \mbox{\lstinline+eq_state_ext s' (update s q x v)+}
    \end{array}
}
\]

\todo{as for choreographies [cit itp], We do this in two steps... first rich labels then normal labels}

\begin{lstlisting}
Inductive SP_To (Defs : DefSetB) :
  Network -> State -> RichLabel -> Network -> State -> Prop :=
  | S_Com N p e B q x B' N' s s' : let v := (eval_on_state e s p) in
    N p = (q!e; B) -> N q = (p?x; B') ->  N' == N ~~ p ~~ q | p[B] | q[B'] ->
    eq_state_ext s' (update s q x v) -> SP_To Defs N s (R_Com p v q x) N' s'
  | S_LSel N p B q Bl Br N' s s' :
    N p = (q(+)left; B) -> N q = (p & Some Bl // Br) ->
    N' == N ~~ p ~~ q | p[B] | q[Bl] -> eq_state_ext s s' ->
    SP_To Defs N s (R_Sel p q left) N' s'
  | S_Then N p b B1 B2 N' s s' :
    N p = (If b Then B1 Else B2) -> beval_on_state b s p = true ->
    N' == N ~~ p | p[B1] -> eq_state_ext s s' -> SP_To Defs N s (R_Cond p) N' s'
  | S_Call N p X N' s s' :
    N p = Call X -> N' == N ~~ p | p[Defs X] -> eq_state_ext s s' ->
    SP_To Defs N s (R_Call X p) N' s'                                     (...)
\end{lstlisting}
As for CC, this relation is then lifted to programs and closed under reflexivity and transitivity, with similar notations.

Transitions are compatible with network equality and state equivalence.
\begin{lstlisting}
Lemma SP_To_eq : forall Defs N tl s1 N' s2 s1' s2',
  eq_state_ext s1 s1' -> eq_state_ext s2 s2' ->
  SP_To Defs N s1 tl N' s2 -> SP_To Defs N s1' tl N' s2'.
Lemma SP_To_Network_eq : forall N1 N1' N2 SPDefs s s' t,
    (N1 == N2) -> SP_To SPDefs N1 s t N1' s' -> SP_To SPDefs N2 s t N1' s'.
Lemma SP_To_Defs_wd : forall Defs Defs', (forall X, Defs X = Defs' X) ->
  forall N s tl N' s', SP_To Defs N s tl N' s' -> SP_To Defs' N s tl N' s'.
\end{lstlisting}
These results are instrumental in some of the later proofs.
The first two are also proven for programs and many-step transitions, while the last one has the
following counterpart.
\begin{lstlisting}
Lemma SPP_To_Defs_stable : forall Defs' N N' tl s s',
  {|Defs N,s|} --[tl]--> {|Defs' N',s'|} -> Defs = Defs'.
\end{lstlisting}

We also prove that transitions are completely determined by the label.
\begin{lstlisting}
Lemma SP_To_deterministic_1 : forall N N1 N2 tl s s1 s2,
  SP_To Defs N s tl N1 s1 -> SP_To Defs N s tl N2 s2 -> N1 == N2.
Lemma SP_To_deterministic_2 : forall N N1 N2 tl s s1 s2,
  SP_To Defs N s tl N1 s1 -> SP_To Defs N s tl N2 s2 -> eq_state_ext s1 s2.
\end{lstlisting}
Finally, we show that the semantics of SP is confluent.
Although this is not required for our main theorem, it is a nice result that confirms our
expectations.

The formalisation of SP consists of 37 definitions, 80 lemmas, and 2000 lines.
% of Coq code.

\section{Merging}
\label{sec:merging}

Intuitively, process implementations can be generated from choreographies recursively, by projecting
each action in the choreography to the corresponding process action -- for example, a value
communication \lstinline+p#e --> qdollarx+ should be projected as a send action \lstinline+q!e+ at
\lstinline+p+, as a receive action \lstinline+p?x+ at \lstinline+q+, and not projected to any other
processes.
%
However, this causes a problem with conditionals.
Projecting a choreography \lstinline+If p ?? b Then Ct Else Ce+ for any process other than
\lstinline+p+, say \lstinline+q+, requires combining the projections obtained for \lstinline+Ct+ and \lstinline+Ce+, such that \lstinline+q+ can ``react'' to whichever choice \lstinline+p+ will make.
This combination is called \emph{merging}~\cite{CHY12}.

Merge is typically defined as a partial function mapping pairs of behaviours to behaviours that
returns a behaviour combining all possible executions of the two input behaviours (if possible).
% An analysis of the semantics of SP quickly shows that this requires the two input behaviours to be
% structurally similar, with the possible exception of branching terms:
For SP, two behaviours can be merged if they are structurally similar, with the possible exception of branching terms: we can merge branchings that offer options on distinct labels.
For example, merging \lstinline+p & (Some B) // None+ with \lstinline+p & None // (Some B')+ yields \lstinline+p & (Some B) // (Some B')+,
allowing \lstinline+If p??b Then (p-->q[left];; q.e --> p;; End) Else (p-->q[right];; End)+ to be projected for \lstinline+q+ as \lstinline+p & Some (p!e; End) // Some End+.
% 
% the merging of a branching term offering a behaviour \lstinline+B+ for label \lstinline+l+ with another branching term offering  one branching offers a behaviour \lstinline+B+ for label \lstinline+l+ and the other does not, then they can be merged by
% keeping the behaviour \lstinline+B+ for \lstinline+l+.

Formalising merging poses the problem that all functions in Coq are total.
Since the type of arguments to branching terms is \lstinline+option Behaviour+,
trying to define merging with type \lstinline+Behaviour -> Behaviour -> option Behaviour+ fails.\footnote{Essentially because it is not possible to distinguish if the behaviour assigned to a
  label is \lstinline+None+ because it was not defined, or because a recursive call to merge
  failed.}

Instead, we define a type \lstinline+XBehaviour+ of extended behaviours, defined as \lstinline+Behaviour+ with an extra constructor \lstinline+XUndefined : XBehaviour+.
Thus, an \lstinline+XBehaviour+ is a \lstinline+Behaviour+ that may contain \lstinline+XUndefined+
subterms.

The connection between \lstinline+Behaviour+ and \lstinline+XBehaviour+ is established by means of
two functions: \lstinline+inject : Behaviour -> XBehaviour+, which isomorphically injects each
constructor of \lstinline+Behaviour+ into the corresponding one in \lstinline+XBehaviour+, and
\lstinline+collapse : XBehaviour -> XBehaviour+ that maps all \lstinline+XBehaviour+s with
\lstinline+XUndefined+ as a subterm to \lstinline+XUndefined+.
The most relevant properties of these functions are:
\begin{lstlisting}
Lemma inject_elim : forall B, exists B', inject B = B' /\ B' <> XUndefined.
Lemma collapse_inject : forall B, collapse (inject B) = inject B.
Lemma collapse_char'' : forall B, collapse B = XUndefined -> forall B', B <> inject B'.
Lemma collapse_exists : forall B, collapse B <> XUndefined -> exists B', B = inject B'.
\end{lstlisting}

Using this type, we first define \lstinline+XMerge+ on \lstinline+XBehaviour+s as follows, where
we report only representative cases.\todo{=? for pid\_dec and expr\_dec with disclaimer}
\begin{lstlisting}
Fixpoint Xmerge (B1 B2:XBehaviour) : XBehaviour := match B1, B2 with
| XEnd, XEnd => XEnd
| XSend p e B, XSend p' e' B' => if Pid_dec p p' && Expr_dec e e'
    then match Xmerge B B' with XUndefined => XUndefined
                              | _ => XSend p e (Xmerge B B') end
    else XUndefined
| XBranching p Bl Br, XBranching p' Bl' Br' => if Pid_dec p p'
    then let BL := match Bl with None => Bl' | Some B =>
        match Bl' with None => Bl | Some B' => Some (Xmerge B B') end end
      in let BR := match Br with None => Br' | Some B =>
        match Br' with None => Br | Some B' => Some (Xmerge B B') end end
      in match BL, BR with Some XUndefined, _ => XUndefined
         | _, Some XUndefined => XUndefined | _, _ => XBranching p BL BR end
    else XUndefined
| XCond e B1 B2, XCond e' B1' B2' => if BExpr_dec e e'
    then match Xmerge B1 B1', Xmerge B2 B2' with XUndefined, _ => XUndefined
         | _, XUndefined => XUndefined | Bt, Be => XCond e Bt Be end
    else XUndefined                                                        (...)
\end{lstlisting}

Using \lstinline+XMerge+ we can straightforwardly define merging.
\begin{lstlisting}
Definition merge B1 B2 := Xmerge (inject B1) (inject B2).
\end{lstlisting}

Function \lstinline+merge+ is shown to be idempotent and commutative.
It is also associative, but this has not been proven: it is not relevant for the remaining results,
and its proof illustrates the major challenge experienced in this stage of the formalisation: all
proofs are by induction on behaviours, and the number of cases generated quickly becomes
unmanageable.
Automation only works to a limited extent, and therefore the proof of associativity was dropped.
(This proof generated $84$ subcases that were not automatically proven. These had to be divided in
further subcases, which could then be dealt with relatively simply. The final case, when all
behaviours are branching terms, then must be divided in the $64$ possible combinations of
defined/undefined branches, which although similar are subtly different. Three of these cases were
proven, leading to our claim that the proof can be completed given sufficient patience.)

The largest set of lemmas about \lstinline+merge+ deals with inversion
results, such as:
%\todo{in the second one, we can replace the part about the right branches by ``\ldots''}
\begin{lstlisting}
Lemma merge_inv_Send : forall B B' p e X, merge B B' = XSend p e X ->
  exists B1 B1', B = p ! e; B1 /\ B' = p ! e; B1' /\ merge B1 B1' = X.
Lemma merge_inv_Branching : forall B B' p Bl Br, merge B B' = XBranching p Bl Br ->
  exists Bl' Bl'' Br' Br'', B = p & Bl' // Br' /\ B' = p & Bl'' // Br''
  /\ (Bl = None -> Bl' = None /\ Bl'' = None) /\ (...)
  /\ (forall BL, Bl = Some BL ->
         (Bl' = None -> exists BL'', Bl'' = Some BL'' /\ BL = inject BL'')
      /\ (Bl'' = None -> exists BL', Bl' = Some BL' /\ BL = inject BL')
      /\ (forall BL' BL'', Bl' = Some BL' /\ Bl'' = Some BL'' -> merge BL' BL'' = BL))
\end{lstlisting}
The omitted parts of lemma \lstinline+merge_inv_Branching+ state properties of \lstinline+Br+ analogous to those of \lstinline+Bl+.
\begin{comment}
Lemma merge_inv_Branching : forall B B' p Bl Br, merge B B' = XBranching p Bl Br ->
  exists Bl' Bl'' Br' Br'', B = p & Bl' // Br' /\ B' = p & Bl'' // Br''
  /\ (Bl = None -> Bl' = None /\ Bl'' = None)
  /\ (Br = None -> Br' = None /\ Br'' = None)
  /\ (forall BL, Bl = Some BL ->
         (Bl' = None -> exists BL'', Bl'' = Some BL'' /\ BL = inject BL'')
      /\ (Bl'' = None -> exists BL', Bl' = Some BL' /\ BL = inject BL')
      /\ (forall BL' BL'', Bl' = Some BL' /\ Bl'' = Some BL'' -> merge BL' BL'' = BL))
  /\ (forall BR, Br = Some BR ->
         (Br' = None -> exists BR'', Br'' = Some BR'' /\ BR = inject BR'')
      /\ (Br'' = None -> exists BR', Br' = Some BR' /\ BR = inject BR')
      /\ (forall BR' BR'', Br' = Some BR' /\ Br'' = Some BR'' -> merge BR' BR'' = BR)).

Lemma merge_inv_Cond : forall B B' b Be Bt, merge B B' = XCond b Be Bt ->
  exists Be' Be'' Bt' Bt'', B = Cond b Be' Bt' /\ B' = Cond b Be'' Bt''
    /\ merge Be' Be'' = Be /\ merge Bt' Bt'' = Bt.
\end{lstlisting}
\end{comment}
Although these proofs amount mostly to induction and case analysis, they suffer from the problems
described for the proof of associativity (thankfully, not to such a dramatic level).
Automation works better here, and the effect of the large number of
subcases is mostly felt in the time required by the \lstinline+auto+ tactic.
Still, the formalisation of merging consists of 6 definitions, 42 lemmas, and 1950 lines %of Coq code
-- giving an average proof length of over 40 lines.

\section{EndPoint Projection}
\label{sec:epp}

We are now ready to define EndPoint Projection (EPP): a partial function that maps programs in CC to
programs in SP. %, provided that the choreography specifies a realisable protocol.
The target instance of SP has the same parameters as CC, except for the set of procedure names:
this is defined as \lstinline+RecVar*Pid+ -- we project an implementation of each procedure from the point of view of each process involved in it.

EPP is partial because of the standard problem of choreography realisability~\cite{BBO12}. Consider the choreography
\lstinline+If p ?? b Then (q#e --> rdollarx) Else End+.
This cannot be implemented without additional communications between \lstinline+p+,
\lstinline+q+ and \lstinline+r+, since the latter processes need to know the result of the
conditional to decide whether to communicate.
We say that this choreography is not \emph{projectable}~\cite{CM20}.

We define EPP in several layers.
First, we define a function \lstinline+bproj : DefSet ->+ \lstinline+Choreography -> Pid -> XBehaviour+ for projecting the behaviour of a single process.
%\lstinline+bproj : DefSet -> Choreography -> Pid -> XBehaviour+.
Intuitively, \lstinline+bproj Defs C p+ attempts to construct \lstinline+p+'s behaviour as specified
by \lstinline+C+; the parameter \lstinline+Defs+ is used to project procedure calls, which depend on
whether \lstinline+p+ participates in the procedure.
Returning an \lstinline+XBehaviour+ instead of an \lstinline+option Behaviour+ has the
advantage of giving information about where exactly merging fails (the location of
\lstinline+XUndefined+ subterms).
This information can be used for debugging, providing information to programmers, or for automatic repair of choreographies~\cite{BB16,CM20}, which we plan to exploit in future work.

We show some illustrative cases of the definition of \lstinline+bproj+.
\begin{lstlisting}
Fixpoint bproj (Defs:DefSet) (C:Choreography) (r:Pid) : XBehaviour := match C with
| p#e --> qdollarx;; C' =>
  if Pid_dec p r then XSend q e (bproj Defs C' r)
                else if Pid_dec q r then XRecv p x (bproj Defs C' r)
                                    else bproj Defs C' r
| p --> q[left];; C' =>
  if Pid_dec p r then XSel q left (bproj Defs C' r)
                 else if Pid_dec q r then XBranching p (Some (bproj Defs C' r)) None
                                     else bproj Defs C' r
| If p ?? b Then C1 Else C2 =>
  if Pid_dec p r then XCond b (bproj Defs C1 r) (bproj Defs C2 r)
                 else Xmerge (bproj Defs C1 r) (bproj Defs C2 r)
| CCBase.Call X => if In_dec P.eq_dec r (fst (Defs X)) then XCall (X,r) else XEnd
(...)
\end{lstlisting}

The next step is generating projections for all relevant processes.
We take the set of processes as a parameter, and collapse all individual projections.
\begin{lstlisting}
Definition epp_list (Defs:DefSet) (C:Choreography) (ps:list Pid)
	: list (Pid * XBehaviour) := map (fun p => (p, collapse (bproj Defs C p))) ps.
\end{lstlisting}

A choreography \lstinline+C+ is projectable wrt \lstinline+Defs+ and \lstinline+ps+ iff
\lstinline+epp_list Defs C ps+ does not contain \lstinline+XUndefined+.
Set \lstinline+Defs:DefSet+ is projectable wrt a set of procedure names \lstinline+Xs+ if
\lstinline+snd (Defs X)+ is projectable wrt \lstinline+Defs+ and \lstinline+fst (Defs X)+ for each
\lstinline+X+ in \lstinline+Xs+.
Projectability of programs is a bit more involved, and we present its Coq formalisation before discussing it.
\begin{lstlisting}
Definition projectable Xs ps P :=
  projectable_C (Procedures P) ps (Main P) /\ projectable_D Xs (Procedures P) /\
  (forall p, In p (CCC_pn (Main P) (fun _ => nil)) -> In p ps) /\
  (forall p X, In X Xs -> In p (fst (Procedures P X)) -> In p ps) /\
  (forall p X, In X Xs -> In p (CCC_pn (snd (Procedures P X)) (fun _ => nil)) -> In p ps).
\end{lstlisting}
The first two conditions simply state that \lstinline+Main P+ and \lstinline+Procedures P+ are
projectable wrt the appropriate parameters.
The remaining conditions state that the sets \lstinline+ps+ and \lstinline+Xs+ include all processes
used in \lstinline+P+ and all procedures needed to execute \lstinline+P+.
(These sets are not necessarily computable, since \lstinline+Xs+ is not required to be finite.
However, in practice, these parameters are known -- so it makes sense to include them in the
definition.)
Function \lstinline+CCC_pn+ returs the set of processes used in a choreography, given the sets of
processes each procedure call is supposed to use.

We now define \lstinline+epp_C+, for compiling projectable choreographies to networks, \lstinline+epp_D+, for compiling projectable sets of procedure definitions, and \lstinline+epp+, for compiling projectable programs.
As these definitions depend on proof terms whose structure needs to be explored, they are done
interactively; afterwards, we prove that these definitions are independent of the proof terms, and otherwise work as expected.
We give a few examples.

\begin{lstlisting}
Lemma epp_C_wd : forall Defs C ps H H', (epp_C Defs ps C H) == (epp_C Defs ps C H').
Lemma epp_C_Com_p : forall Defs ps C p e q x HC HC', In p ps ->
  epp_C Defs ps (p#e-->qdollarx;;C) HC p = q!e; epp_C Defs ps C HC' p.
Lemma epp_C_Cond_r : forall Defs ps p b C1 C2 HC HC1 HC2 r, p <> r ->
  inject (epp_C Defs ps (If p ?? b Then C1 Else C2) HC r)
  = merge (epp_C Defs ps C1 HC1 r) (epp_C Defs ps C2 HC2 r).
\end{lstlisting}
\begin{comment}
Lemma epp_C_out : forall Defs C ps H p, ~In p ps -> epp_C Defs ps C H p = End.
Lemma epp_C_char : forall Xs ps Defs C HP HC,
  (Net (epp Xs ps {| Procedures := Defs; Main := C |} HP)) == (epp_C Defs ps C HC).
Lemma epp_C_Cond_p : forall Defs ps p b C1 C2 HC HC1 HC2, In p ps ->
  epp_C Defs ps (If p ?? b Then C1 Else C2) HC p
  = If b Then (epp_C Defs ps C1 HC1 p) Else (epp_C Defs ps C2 HC2 p).
\end{lstlisting}
\end{comment}

Furthermore, we need a number of inversion lemmas for branching and conditionals that restrict
possible projections of subterms of a choreography.
\begin{comment}
Lemma epp_C_not_Branching_None_None : forall Defs ps C HC p q,
  epp_C Defs ps C HC p <> q & None // None.
\end{comment}
\begin{lstlisting}
Lemma epp_C_Sel_Branching_l : forall Defs ps C HC p q Bp Bl Br,
  epp_C Defs ps C HC p = q (+) left; Bp ->
  epp_C Defs ps C HC q = p & Bl // Br -> Bl <> None /\ Br = None.
Lemma epp_C_Cond_Send_inv : forall Defs ps p b C1 C2 HC HC1 HC2 r q e B,
  epp_C Defs ps (If p ?? b Then C1 Else C2) HC r = q ! e; B ->
  exists B1 B2, epp_C Defs ps C1 HC1 r = q ! e; B1
  /\ epp_C Defs ps C2 HC2 r = q ! e; B2 /\ merge B1 B2 = inject B.
\end{lstlisting}

As defined so far, projectability of \lstinline+C+ does not imply projectability of choreographies
that \lstinline+C+ can transition to.
This is due to the way runtime terms are projected: \lstinline+RT_Call X ps C'+ is projected as a
call to \lstinline+(X,p)+ if \lstinline+p+ is in \lstinline+ps+, and as the projection of
\lstinline+C'+ otherwise.
Our definition of projectability allows in principle for \lstinline+C+ to be unprojectable for a
process in \lstinline+ps+, which would make it unprojectable after transition.
% 
That this situation does not arise is a consequence of the intended usage of runtime terms:
initially \lstinline+C'+ is obtained from the definition of a procedure, and \lstinline+ps+ is the
set of processes used in this procedure.
Afterwards \lstinline+ps+ only shrinks, while \lstinline+C'+ may change due to execution of actions
outside \lstinline+ps+.
We capture these conditions in the notion of strong projectability, whose
representative case is:
\begin{lstlisting}
Fixpoint strongly_projectable Defs (C:Choreography) (r:Pid) : Prop := 
match C with | RT_Call X ps C => strongly_projectable Defs C r /\
     (forall p, In p ps -> In p (fst (Defs X))
          /\ Xmore_branches (bproj Defs (snd (Defs X)) p) (bproj Defs C p))      (...)
\end{lstlisting}
The relation \lstinline+Xmore_branches+ is explained in the next section: it is a semantic
characterisation of how the projection of \lstinline+bproj Defs snd (Defs X) p+ may change due to
execution of actions not involving \lstinline+p+ in \lstinline+snd (Defs X)+.

% By definition, 
Projectability and strong projectability coincide for initial choreographies.
Furthermore, we state and prove lemmas that show that strong projectability of \lstinline+C+ imply
strong projectability of any choreography that \lstinline+C+ can transition to.
%% -- excluding procedure calls, since those are covered by the definition of projectable program.
\begin{comment}
\begin{lstlisting}
Lemma projectable_inv_Com : forall Xs ps Defs p e q x C,
  projectable Xs ps (CCBase.Build_Program Defs (p#e-->qdollarx;;C)) ->
  projectable Xs ps (CCBase.Build_Program Defs C).
Lemma projectable_inv_RT_Call : forall Xs ps Defs X p ps' C,
  projectable Xs ps (CCBase.Build_Program Defs (RT_Call X ps' C)) ->
  collapse (bproj Defs C p) <> XUndefined ->
  projectable Xs ps (CCBase.Build_Program Defs (RT_Call X (set_remove_pid p ps') C)).
\end{lstlisting}
\end{comment}

\section{Pruning}
\label{sec:pruning}

The key ingredient for our correspondence result is a relation on behaviours usually called \emph{pruning} \cite{CHY12,CM13}.
% , which we briefly mentioned at the end of the previous section.
Pruning relates two behaviours that differ only in that one offers more options in branching terms
than the other; we formalise this relation with the more suggestive name of
\lstinline+more_branches+ (in line with~\cite{M21itc}), and we include some illustrative cases of its definition.
\begin{lstlisting}
Inductive more_branches : Behaviour -> Behaviour -> Prop :=
| MB_Send p e B B': more_branches B B' -> more_branches (p ! e; B) (p ! e; B')
| MB_Branching_None_None p mBl mBr :
    more_branches (p & mBl // mBr) (p & None // None)
| MB_Branching_Some_Some p Bl Bl' Br Br' :
    more_branches Bl Bl' -> more_branches Br Br' ->
    more_branches (p & Some Bl // Some Br) (p & Some Bl' // Some Br')           (...)
\end{lstlisting}
The need for pruning arises naturally when one considers what happens when a choreography executes a
conditional at a process \lstinline+p+.
In the continuation, only one of the branches will be kept.
However, no process other than \lstinline+p+ knows that this action has been executed; therefore, in the
choreography's projection both behaviours will still be available for all other processes.
Given the definition of merging, this means that these processes' behaviours may contain more
branches than those of the projection of the choreography after reduction.

Pruning is naturally related to merging, as stated in the following lemmas.
\begin{lstlisting}
Lemma more_branches_char : forall B B', more_branches B B' <-> merge B B' = inject B.
Lemma merge_more_branches :
  forall B1 B2 B, merge B1 B2 = inject B -> more_branches B B1.
\end{lstlisting}

Pruning is also reflexive and transitive.
Further, if two behaviours have an upper bound according to pruning, then their merge is defined and is their lub.
\begin{lstlisting}
Lemma more_branches_merge : forall B B1 B2, more_branches B B1 ->
  more_branches B B2 -> exists B', merge B1 B2 = inject B' /\ more_branches B B'.
\end{lstlisting}

Finally, two behaviours that are lower than two mergeable behaviours are themselves mergeable.
\begin{lstlisting}
Lemma more_branches_merge_extend : forall B1 B2 B1' B2' B,
  more_branches B1 B1' -> more_branches B2 B2' -> merge B1 B2 = inject B ->
  exists B', merge B1' B2' = inject B' /\ more_branches B B'.
\end{lstlisting}

These two results are key ingredients to the cases dealing with conditionals in the proof of the EPP
Theorem.
% As before, t
They require extensive case analysis (512 cases for the last lemma, of which 81 are not
automatically solved by Coq's \lstinline+inversion+ tactic even though most are contradictory).
Analogous versions of some of these lemmas also need to be extended to \lstinline+XBehaviour+s,
which is straightforward.

Pruning extends pointwise to networks, which we denote as \lstinline+N >> N'+.
The key result is that, due to how the semantics of SP is defined, pruning a network cannot add new
transitions.
\begin{lstlisting}
Lemma SP_To_more_branches_N : forall Defs N1 s N2 s' Defs' N1' tl,
  SP_To Defs N1 s tl N2 s' -> N1' >> N1 -> (forall X, Defs X = Defs' X) ->
  exists N2', SP_To Defs' N1' s tl N2' s' /\ N2' >> N2.
\end{lstlisting}
The reciprocal of this result only holds for choreography projections, and is proven after the
definition of EPP.

The formalisation of pruning includes 3 definitions, 25 lemmas, and 950 lines.
% of Coq code.
This file is imported in the file containing the definition of EPP (previous section), but we
delayed its presentation as its motivation is clearer after seeing those definitions.

\todo{example of process projections, example of merging something from the example}

\todo{somewhere example saying that the EPP of what we had is exactly what we wrote in the section on processes}

\section{EPP Theorem}
\label{sec:epp-theorem}

We now prove the operational correspondence between choreographies and their projections, in two directions: if a choreography can make a transition, then its projection can make the same transition; and if the
projection of a choreography can make a transition, then so can the choreography.
The results of the transitions are not directly related by projection, since choreography transitions
may eliminate some branches in the projection; thus, establishing the correspondence for multi-step transitions
requires some additional lemmas on pruning.

\paragraph{Preliminaries.}
Both directions of the correspondence depend on a number of results relating choreography
transitions and their projections.
These results follow a pattern: the results for communications state a precise
correspondence; the ones for conditionals include pruning in their conclusions; and the ones for
procedure calls require additional hypotheses on the set of procedure definitions.
\begin{lstlisting}
Lemma CCC_To_bproj_Sel_p : forall Defs C s C' s' p q l,
  strongly_projectable Defs C p ->
  CCC_To Defs C s (CCBase.TL.R_Sel p q l) C' s' ->
  exists Bp, bproj Defs C p = XSel q l Bp /\ bproj Defs C' p = Bp.

Lemma CCC_To_bproj_Call_p : forall Defs C s C' s' p X Xs,
  strongly_projectable Defs C p ->
  (forall Y, In Y Xs -> strongly_projectable Defs (snd (Defs Y)) p) ->
  (forall Y, set_incl_pid (CCC_pn (snd (Defs Y)) (fun X => fst (Defs X)))
                          (fst (Defs Y))) ->
  In X Xs ->  CCC_To Defs C s (CCBase.TL.R_Call X p) C' s' ->
  bproj Defs C p = XCall (X,p)
  /\ Xmore_branches (bproj Defs (snd (Defs X)) p) (bproj Defs C' p).
\end{lstlisting}
These lemmas are simple to prove by induction on \lstinline+C+.
The tricky part is getting the hypotheses strong enough that the thesis holds, and weak enough that any well-formed program will satisfy them throughout its entire execution.

From these results, it follows that projectability is preserved by choreography reductions.
This property is needed even to state the EPP Theorem, since we can only compute projections of projectable programs.
\begin{lstlisting}
Lemma CCC_To_projectable : forall P Xs ps,
  Program_WF Xs P -> well_ann P -> projectable Xs ps P ->
  (forall p, In p ps -> strongly_projectable (Procedures P) (Main P) p) ->
  (forall p, In p (CCC_pn (Main P) (Vars P)) -> In p ps) ->
  (forall p X, In X Xs -> In p (Vars P X) -> In p ps) ->
  forall s tl P' s', (P,s) --[tl]--> (P',s') -> projectable Xs ps P'.
\end{lstlisting}
Some of the hypotheses from the previous lemmas are encapsulated in the first two conditions:
well-formedness of \lstinline+P+, which ensures that any runtime term \lstinline+RT_Call X qs C+ in
\lstinline+Main P+ only includes processes in \lstinline+qs+ that are declared to be used by
\lstinline+X+ (this trivially holds if \lstinline+Main P+ is initial, and is preserved throughout
execution); and well-annotation of \lstinline+P+, i.e., the processes used in any
procedure are a subset of those it declares.\footnote{Equality is not necessary, and it would make
  this property harder to prove.}
The remaining hypotheses state, as before, that
% the sets
\lstinline+Xs+ and \lstinline+ps+ include
all processes and procedures relevant for executing \lstinline+Main P+.

Similarly, we prove that strong projectability is preserved by transitions.

\paragraph{Completeness.}
In the literature, completeness of EPP is proven by structural induction on the derivation of the transition performed by the choreography $C$.
For each case, we look at how a transition for $C$ can be derived,
% (typically in two ways --
% consuming the first action in $C$ or delaying it),
and show that the projection of $C$ can make the
same transition to a network with more branches than the projection of $C'$.
The proof is lengthy, but poses no surprises.
\begin{lstlisting}
Lemma EPP_Complete : forall P Xs ps, Program_WF Xs P -> well_ann P -> forall HP,
  (forall p, In p ps -> strongly_projectable (Procedures P) (Main P) p) ->
  (forall p, In p (CCC_pn (Main P) (Vars P)) -> In p ps) ->
  (forall p X, In X Xs -> In p (Vars P X) -> In p ps) ->
  forall s tl P' s', (P,s) --[tl]--> (P',s') ->
  exists N tl', (epp Xs ps P HP,s) --[tl']--> (N,s')
    /\ Procs N = Procs (epp Xs ps P HP) /\ forall H, Net N >> Net (epp Xs ps P' H).
\end{lstlisting}
By combining with the earlier results on pruning, we immediately obtain the generalisation for
multi-step transitions.

\paragraph{Soundness.}
% In the literature, soundness is proven
We prove soundness by case analysis on the transition made by the network, and
then by induction on the choreography inside each case.
For convenience, we split this proof in separate proofs, one for each transition.
Again we observe that the results for communications are stronger than those for conditionals, which
in turn have fewer hypotheses than those for procedure calls.
\begin{lstlisting}
Lemma SP_To_bproj_Com : forall Defs Defs' ps C HC s N' s' p x q v,
  (forall p, In p ps -> strongly_projectable Defs C p) ->
  (forall p, In p (CCC_pn C (fun X => fst (Defs X))) -> In p ps) ->
  SP_To Defs' (epp_C Defs ps C HC) s (R_Com p v q x) N' s' ->
  exists C', CCC_To Defs C s (CCBase.TL.R_Com p v q x) C' s'
  /\ forall HC', N' == (epp_C Defs ps C' HC').

Lemma SP_To_bproj_Call : forall Defs Defs' ps C HC s N' s' p X Xs,
  Choreography_WF C -> within_Xs Xs C -> In X Xs ->
  (forall p, In p ps -> strongly_projectable Defs C p) ->
  (forall p, In p (CCC_pn C (fun X => fst (Defs X))) -> In p ps) ->
  (forall p HX, Defs' (X,p) = epp_C Defs ps (snd (Defs X)) HX p) ->
  (forall p X, In p (CCC_pn (snd (Defs X)) (fun Y => fst (Defs Y))) ->
          In p (fst (Defs X))) ->
  (forall X, In X Xs -> projectable_C Defs ps (snd (Defs X))) ->
  (forall X, In X Xs -> initial (snd (Defs X))
    /\ forall p, In p (fst (Defs X)) -> In p ps) ->
  SP_To Defs' (epp_C Defs ps C HC) s (R_Call (X,p) p) N' s' ->
  exists C', CCC_To Defs C s (CCBase.TL.R_Call X p) C' s'
  /\ forall HC', N' >> epp_C Defs ps C' HC'.
\end{lstlisting}
By contrast with completeness, all these lemmas are complex to prove: each requires around
300 lines of Coq code.
The proofs are similar,
% making it easy to adapt previous ones,
but still different enough that the
task is not completely mechanic.

The last ingredient is a lemma of practical interest on procedure names: each process only uses ``its'' copy of the original procedure names.
\begin{lstlisting}
Lemma SP_To_bproj_Call_name : forall Defs Defs' ps C HC s N' s' p X,
  SP_To Defs' (epp_C Defs ps C HC) s (R_Call X p) N' s' ->
  exists Y, X = (Y,p) /\ X_Free Y C.
\end{lstlisting}
This lemma is not only crucial in the proof of the next theorem, but also interesting in itself: it
shows that the set of procedure definitions can be fully distributed among the processes with no
duplications.

\begin{lstlisting}
Lemma EPP_Sound : forall P Xs ps, Program_WF Xs P -> well_ann P -> forall HP,
  (forall p, In p ps -> strongly_projectable (Procedures P) (Main P) p) ->
  (forall p, In p (CCC_pn (Main P) (Vars P)) -> In p ps) ->
  (forall p X, In X Xs -> In p (Vars P X) -> In p ps) ->
  forall s tl N' s', (epp Xs ps P HP,s) --[tl]-->(N',s') ->
  exists P' tl', (P,s) --[tl']--> (P',s') /\ forall H, Net N' >> Net (epp Xs ps P' H).
\end{lstlisting}

Generalising this result to multi-step transitions requires showing that pruning does not eliminate
possible transitions of a network.
This is in general not true, but it holds when the pruned network is the projection of a
choreography.

\begin{lstlisting}
Lemma SP_To_more_branches_N_epp : forall Defs N1 s N2 s' tl Defs' ps C HC,
  N1 >> epp_C Defs' ps C HC -> SP_To Defs N1 s tl N2 s' ->
  exists N2', SP_To Defs (epp_C Defs' ps C HC) s tl N2' s' /\ N2 >> N2'.
\end{lstlisting}

The formalisation of EPP and the proof of the EPP theorem consists of 13 definitions, 110 lemmas, and
4960 lines of Coq code.
The proof of the EPP Theorem and related lemmas make up for around 75\%\ of this size.

\section{Discussion and Conclusion}

We have successfully formalised a translation from a Turing complete choreographic language into a process calculus and proven its correctness in terms of an operational correspondence result.
This formalisation showed that the proof techniques used in the literature are correct, and
identified only missing minor assumptions about runtime terms that trivially hold when these are
used as intended.
To the best of our knowledge, this is the first time such a correspondence has been formalised for a
full-fledged (Turing complete) choreographic language.

The complexity of the formalisation, combined with the similarities between several of the proofs,
means that future extensions would benefit from exploiting semi-automatic generation of proof
scripts.

Combining these results with those from~\cite{CMP21} would yield a proof that SP is also
Turing complete.
Unfortunately, the choreographies used in the proof of Turing completeness in~\cite{CMP21} are not projectable, but they can be made so automatically, by means of an amendment (repair) procedure~\cite{CM20}.
In future work, we plan to formalise amendment in order to obtain this result.
        
\bibliographystyle{splncs04}
\bibliography{biblio}
\end{document}
